\documentclass[11pt,a4paper]{article}
\usepackage[T1]{fontenc}
\usepackage[utf8]{inputenc}
\usepackage[a4paper,margin=2.5cm]{geometry}
\usepackage{xcolor}
\usepackage{hyperref}
\usepackage{enumitem}
\usepackage{booktabs}
\usepackage{titlesec}
\usepackage{mdframed}
\usepackage{parskip}
\usepackage{fancyhdr}
\definecolor{navyblue}{RGB}{31,56,100}
\definecolor{medblue}{RGB}{46,117,182}
\definecolor{lightblue}{RGB}{68,114,196}
\hypersetup{colorlinks=true,linkcolor=medblue,urlcolor=medblue}
\titleformat{\section}{\large\bfseries\color{navyblue}}{}{0em}{\thesection\quad}
\titleformat{\subsection}{\normalsize\bfseries\color{medblue}}{}{0em}{\thesubsection\quad}
\titleformat{\subsubsection}{\normalsize\bfseries\color{lightblue}}{}{0em}{\thesubsubsection\quad}
\pagestyle{fancy}
\fancyhf{}
\lhead{\color{navyblue}\small\leftmark}
\rfoot{\small\thepage}
\mdfdefinestyle{warnbox}{backgroundcolor=yellow!15,linecolor=orange,linewidth=1pt,
  innertopmargin=6pt,innerbottommargin=6pt,innerleftmargin=8pt,innerrightmargin=8pt}
\mdfdefinestyle{recbox}{linecolor=medblue,linewidth=3pt,topline=false,bottomline=false,rightline=false,
  innerleftmargin=10pt,innerrightmargin=8pt,innertopmargin=4pt,innerbottommargin=4pt}

\begin{document}
\begin{titlepage}
\centering\vspace*{3cm}
{\Huge\bfseries\color{navyblue} Presentation Outline\\[0.5em]}
{\large\color{medblue} Patient Monitoring via Structured Electronic Health Records\\[2em]}
{\normalsize SwissCare Health Network AG --- AI Project Course}\\[0.5em]
{\small\today}
\end{titlepage}
\tableofcontents\newpage
\section{Presentation Outline}
\subsection{Patient Monitoring via Structured Electronic Health Records}
\subsubsection{Preventing Adverse Diabetic Events Through ML Patient Monitoring}
\textbf{Presenting team (consultants):} Sandy, Nicholas, Timarian
\textbf{Client:} SwissCare Health Network AG
\textbf{Duration:} 40 min + Q\&A

\noindent\rule{\linewidth}{0.4pt}

\begin{mdframed}[style=warnbox]
\end{mdframed}
\subsection{\textbf{[!]} TO BE DISCUSSED BEFORE FINALISING}

\subsubsection{Dashboard transition --- when?}
\begin{itemize}
\item \textbf{Option A --- Switch at Section 4 only:} Slides 1--12 as deck, then live dashboard for Results \& Evaluation (Section 4). Cleanest split: deck tells the story, dashboard shows the evidence.
\item \textbf{Option B --- Switch at Section 2:} Present Sections 2, 3, and 4 entirely from the dashboard (one tab each: Data \& Methodology, Architecture, Results). More impressive technically; riskier if the demo has issues; requires dashboard tabs to be presentation-ready.
\item \textbf{Recommendation to discuss:} Option A is safer for a 40-min slot. Option B is more compelling if the dashboard is polished and the tabs map cleanly onto the section structure.
\end{itemize}

\subsubsection{Review checklist --- fill in before building slides}
\begin{itemize}
\item [ ] \textbf{Section 2:} Confirm exact data quality figures (\% duplicates, \% orphaned records, \% missing per feature)
\item [ ] \textbf{Section 2:} Confirm number of clinical note records and format breakdown (free text vs. semi-structured)
\item [ ] \textbf{Section 2:} Confirm preprocessing steps and which cleaning issues were most impactful
\item [ ] \textbf{Section 3:} Confirm pipeline runtime on the hardware you'll demo from (\textasciitilde{}20 min estimate)
\item [ ] \textbf{Section 3:} Confirm key technical challenges --- align with what Timarian's error analysis found
\item [ ] \textbf{Section 4:} Fill in actual PR-AUC and ROC-AUC numbers from the trained XGBoost model
\item [ ] \textbf{Section 4:} Fill in the "100 alerts" figure: X caught / Y false alarms at the operating threshold
\item [ ] \textbf{Section 4:} Insert Timarian's error analysis findings (false negative and false positive profiles)
\item [ ] \textbf{Section 4:} Insert fairness results --- which demographic group shows the largest performance gap
\item [ ] \textbf{Section 5:} Confirm the ROI estimate (cost per prevented readmission $\times$ expected prevention rate)
\item [ ] \textbf{Section 5:} Confirm Alice/Bob patient IDs or synthetic examples for the SHAP force plot demo
\item [ ] \textbf{All sections:} REVIEW SECTIONS 2--6 IN DETAIL before finalising slide content
\end{itemize}

\noindent\rule{\linewidth}{0.4pt}

\subsection{Narrative Thread: Alice and Bob}

Run Alice and Bob as a recurring thread to ground technical content in a human story. Introduce them in the problem statement, return to them at each major section transition.

\textbf{Alice}, 67, Type 2 diabetic in Worcester, MA. HbA1c creeping up for 18 months. Missed her last two check-ups. eGFR declining. Her deterioration is documented only in free-text clinical notes --- inaccessible to any monitoring system. Without intervention, she is 30 days from a hospitalisation.

\textbf{Bob}, 52, Type 2 diabetic in Boston. HbA1c borderline, but he had an ER visit 6 weeks ago and was hospitalised twice in the past year. Care fragmented across three providers. Each sees only part of the picture.

Use them at: Slide 3 (problem), Slide 6 (data motivation), Slide 13 (live SHAP demo), Slide 18 (outcome / close).

\noindent\rule{\linewidth}{0.4pt}

\subsection{Timing Budget (40 min)}

\begin{center}
\begin{tabular}{lll}
\toprule
Section & Slides & Time \\
\midrule
1. Introduction \& Recommendation & 1--5 & \textasciitilde{}8 min \\
2. Data \& Methodology & 6--9 & \textasciitilde{}9 min \\
3. Implementation \& Architecture & 10--12 & \textasciitilde{}6 min \\
4. Results \& Evaluation & 13--17 & \textasciitilde{}11 min $\rightarrow$ \textbf{live dashboard} \\
5. Key Takeaways & 18--19 & \textasciitilde{}6 min \\
\bottomrule
\end{tabular}
\end{center}

\noindent\rule{\linewidth}{0.4pt}

\subsection{Section 1 --- Introduction \& Recommendation (\textasciitilde{}8 min)}
\emph{Pyramid structure: lead with the answer. The rest of the presentation is the support.}

\subsubsection{Slide 1 --- Title \& Team}
\begin{itemize}
\item \textbf{Title:} Patient Monitoring via Structured Electronic Health Records
\item \textbf{Subtitle:} Preventing Adverse Diabetic Events Through ML Patient Monitoring
\item \textbf{Consulting team:} Sandy, Nicholas, Timarian
\item \textbf{Client:} SwissCare Health Network AG
\item \textbf{Date / Engagement:} [date]
\end{itemize}

\noindent\rule{\linewidth}{0.4pt}

\subsubsection{Slide 2 --- Our Recommendation \emph{(the pyramid top)}}
\emph{This is the most important slide. State the conclusion before the audience has to ask for it.}

\begin{mdframed}[style=recbox]
\textbf{"We recommend deploying an XGBoost pipeline with 22 clinically-grounded features to predict 30-day adverse events in diabetic patients. At the chosen operating threshold, the system identifies [X]\% of high-risk patients 48--72 hours before adverse events, with an estimated net ROI of [CHF/USD X] per year based on prevented readmissions. The system carries no material bias risk across demographic subgroups and is fully compliant with EU AI Act high-risk requirements, GDPR/revFADP, and HL7 FHIR standards."}
\end{mdframed}

\begin{itemize}
\item Three supporting pillars --- and we will walk through each:
\begin{enumerate}
\item \textbf{It works:} PR-AUC of [X] vs. best clinical rule baseline of 0.155 --- a [Y]x improvement
\item \textbf{It can be trusted:} SHAP explanations per alert, demographic fairness validated, full compliance architecture documented
\item \textbf{It is worth deploying:} ROI positive from [X] prevented readmissions per year; pipeline runs in \textasciitilde{}20 min on standard hardware; no GPU required
\end{enumerate}
\end{itemize}

\noindent\rule{\linewidth}{0.4pt}

\subsubsection{Slide 3 --- The Problem: Alice and Bob}
\begin{itemize}
\item Two diabetic patients. Two trajectories heading toward a preventable crisis.
\item \textbf{Alice (67, Worcester MA):} 6 unread free-text clinical notes flagging worsening glycemic control. No alert fired. No doctor called. 30 days from hospitalisation.
\item \textbf{Bob (52, Boston):} 3 providers, 2 recent hospitalisations, fragmented data. No single person has the full picture.
\item \textbf{The gap:} 70\% of clinical information at SwissCare lives in unstructured notes --- inaccessible to any systematic monitoring today
\item \textbf{The cost:} Average preventable hospitalisation costs \$15,000--\$25,000; each missed readmission is a patient outcome failure and a balance sheet hit
\item \emph{"Our engagement: read what the system isn't reading, and flag Alice and Bob before it's too late."}
\end{itemize}

\noindent\rule{\linewidth}{0.4pt}

\subsubsection{Slide 4 --- Scope \& Mandate}
\begin{itemize}
\item \textbf{Client:} SwissCare Health Network AG --- hospitals and affiliated clinics, Switzerland \& EU
\item \textbf{Dataset:} 50,000+ diabetic patients, Massachusetts (Synthea-generated EHR data)
\item \textbf{Prediction target:} 30-day high-risk adverse event (ER visit, hospitalisation, acute complication, HbA1c $\geq$ 10\%)
\item \textbf{Engagement deliverables:}
\begin{enumerate}
\item End-to-end ML pipeline: raw EHR $\rightarrow$ structured profiles $\rightarrow$ risk scores $\rightarrow$ explainable alerts
\item Compliance documentation: EU AI Act, GDPR/revFADP Article 22, HL7 FHIR, SNOMED/ICD-10
\item Monitoring dashboard: live risk scoring with SHAP explanations, configurable alert thresholds
\end{enumerate}
\item \textbf{Timeline:} 22-week project
\end{itemize}

\noindent\rule{\linewidth}{0.4pt}

\subsubsection{Slide 5 --- Stakeholders \& Their Success Metrics}
\begin{itemize}
\item \textbf{Chief Medical Officer:} Adverse event prediction rate, false positive rate --- \emph{does the alert tell me something clinically actionable?}
\item \textbf{Emergency Department Director:} Alert recall, lead time --- \emph{does it catch every avoidable ER visit early enough?}
\item \textbf{Nursing Leadership:} False positive rate, alert interpretability --- \emph{is every alert worth my team's time?}
\item \textbf{Compliance Officer:} Zero regulatory violations, documentation completeness --- \emph{can we pass an EU AI Act audit?}
\item \textbf{Chief Information Officer:} System reliability, FHIR compliance rate --- \emph{does it integrate with existing systems?}
\item \textbf{Patients (Alice \& Bob):} Better outcomes, right to understand automated decisions affecting their care
\end{itemize}

\noindent\rule{\linewidth}{0.4pt}

\subsection{Section 2 --- Data \& Methodology (\textasciitilde{}9 min)}
\emph{If Option B (dashboard from Section 2): transition here. Dashboard Tab 1: Data \& Methodology.}

\subsubsection{Slide 6 --- The Data}
\begin{itemize}
\item \textbf{Structured EHR:} 50,000+ patient records across 7 relational tables --- patients, encounters, conditions, medications, observations, procedures, allergies
\item \textbf{Unstructured clinical notes:} \textasciitilde{}25\% of encounters exist only in free-text and semi-structured reports --- this is where Alice's deterioration is hiding
\begin{itemize}
\item [confirm total note count and format split]
\end{itemize}
\item \textbf{Geography:} Massachusetts, US; patient ZIP codes available for socioeconomic analysis
\item \textbf{Target population:} Diabetic patients identified by ICD-10 and SNOMED condition codes
\item \textbf{Prediction label:} Binary --- high-risk adverse event (ER visit, hospitalisation, acute complication, HbA1c $\geq$ 10\%) within the next 30 days
\item \emph{"Alice is one of these 50,000 patients. Her risk lives in the notes."}
\end{itemize}

\noindent\rule{\linewidth}{0.4pt}

\subsubsection{Slide 7 --- Data Quality \& Preprocessing}
\begin{itemize}
\item \textbf{Key quality issues encountered:}
\begin{itemize}
\item Class imbalance: only \textasciitilde{}3\% of patient-day instances are positive --- predicting "low risk" always gives 97\% accuracy and zero clinical value
\item OCR errors in scanned notes: character substitutions corrupt lab values before extraction
\item HTML artifacts: residual tags from web EHR exports must be stripped
\item Missing observations: no recent HbA1c, eGFR, or BP for many patients --- missing data is the norm
\item Duplicate and orphaned records: [X]\% required deduplication; orphaned notes resolved [confirm figure]
\end{itemize}
\item \textbf{Pipeline response:}
\begin{itemize}
\item HTML stripping $\rightarrow$ OCR correction (configurable dictionary, safelist for valid codes like "HbA1c")
\item Regex + NER extraction of vitals, labs, medications with LOINC code assignment
\item Range validation: extracted values outside clinical reference ranges discarded
\item All features computed with strict \texttt{cutoff\_time} --- no future data can contaminate training
\end{itemize}
\end{itemize}

\noindent\rule{\linewidth}{0.4pt}

\subsubsection{Slide 8 --- Feature Engineering: 22 Features, 4 Families}
\begin{itemize}
\item \textbf{Why 22 features?} Clinically motivated --- every feature maps to a known risk indicator; no black-box statistical selection
\item \emph{Temporal counters (5):} days since last HbA1c, encounter, ER visit, hospitalisation, medication change --- capture care engagement
\item \emph{Rolling windows (4):} encounters, ER visits, hospitalisations, medication changes over 90--365 days --- capture trajectory
\item \emph{Clinical values (5):} current HbA1c, systolic BP, diastolic BP, eGFR, BMI category --- the clinical picture today
\item \emph{Trend features (3):} HbA1c, BP, eGFR slope over 6--12 months --- is the patient improving or deteriorating?
\item \emph{Composite risk (5):} active medication count, complication count, care gap count, longest care gap, age --- disease burden summary
\item \textbf{Extension tested:} +3 ZIP-code census features (population density, provider density, insurance coverage) $\rightarrow$ near-zero correlations (r < 0.006) --- not deployed; environmental factors do not predict 30-day clinical outcomes when strong clinical features are present
\end{itemize}

\noindent\rule{\linewidth}{0.4pt}

\subsubsection{Slide 9 --- Model Selection \& Validation}
\begin{itemize}
\item \textbf{Selection rationale:} Healthcare context demands explainability --- clinicians cannot act on black-box scores; GDPR Art. 22 requires meaningful explanations; tree methods are well-suited to tabular features with missing values
\end{itemize}

\begin{center}
\begin{tabular}{llll}
\toprule
Model & PR-AUC & ROC-AUC & Why tested \\
\midrule
Random / Majority class & \textasciitilde{}0.03 / negligible & \textasciitilde{}0.50 & Establishes floor; exposes why accuracy is meaningless \\
HbA1c > 9\% rule & \textasciitilde{}0.155 & --- & Clinical rule benchmark \\
Logistic Regression & [result] & [result] & Linear ceiling; L2 shrinks weak features \\
Random Forest & [result] & [result] & Ensemble variance reduction \\
\textbf{XGBoost + SHAP} & \textbf{[result]} & \textbf{[result]} & \textbf{Best PR-AUC; native NaN; full explainability $\rightarrow$ selected} \\
\bottomrule
\end{tabular}
\end{center}

\begin{itemize}
\item \textbf{Validation approach:} Patient-level split (no patient in both train and test) $\rightarrow$ 5-fold GroupKFold stability $\rightarrow$ temporal validation (train on past patients, test on future) $\rightarrow$ robustness under degraded missing data rates
\item \textbf{Class imbalance:} \texttt{scale\_pos\_weight} in XGBoost; \texttt{eval\_metric='aucpr'} to optimise directly for the clinical metric
\end{itemize}

\noindent\rule{\linewidth}{0.4pt}

\subsection{Section 3 --- Implementation \& Technical Architecture (\textasciitilde{}6 min)}
\emph{If Option B (dashboard from Section 2): Dashboard Tab 2: Architecture.}

\subsubsection{Slide 10 --- Pipeline Architecture}
\emph{Visual: left-to-right flowchart --- clean, consulting-style diagram}

\begin{mdframed}[backgroundcolor=gray!10,linecolor=gray,linewidth=0.5pt]
\begin{verbatim}
[Raw EHR CSVs]     [Clinical Notes]
       ↓                  ↓
     [OCR / HTML Cleaning]
              ↓
   [Regex + NER Extraction → LOINC codes]
              ↓
   [FHIR-structured Patient Profiles]
              ↓
   [22-Feature Computation (cutoff-safe)]
              ↓
        [XGBoost Scorer]
              ↓
     [SHAP Explanation Layer]
              ↓
     [Alert Dashboard (Streamlit)]
\end{verbatim}
\end{mdframed}

\begin{itemize}
\item Modular \texttt{EHRPipeline} class: each stage chains cleanly, designed for both batch and single-patient real-time query
\item Single-patient prediction: sub-second once profiles are loaded
\item Alert levels: CRITICAL $\geq$ 80\%, HIGH $\geq$ 60\%, MODERATE $\geq$ 40\%, LOW < 40\% --- configurable per ward
\end{itemize}

\noindent\rule{\linewidth}{0.4pt}

\subsubsection{Slide 11 --- System Architecture \& Deployment}
\begin{itemize}
\item \textbf{Data sources:} Synthea CSV tables + free-text notes (extensible to live FHIR API in production)
\item \textbf{Core API:} \texttt{predict\_patient(patient\_id, date)} $\rightarrow$ risk probability; \texttt{explain\_patient(patient\_id, date)} $\rightarrow$ SHAP decomposition with clinical language
\item \textbf{Monitoring dashboard:} Streamlit --- patient explorer, model performance, feature importance, per-patient SHAP force plots
\item \textbf{Computational footprint:} Full pipeline \textasciitilde{}20 min on standard laptop / Colab; no GPU required; designed to run within hospital data perimeter --- no patient data transmitted externally
\item \textbf{Integration points:} FHIR-structured outputs, LOINC-coded observations, SNOMED event definitions --- compatible with any FHIR-compliant hospital system
\end{itemize}

\noindent\rule{\linewidth}{0.4pt}

\subsubsection{Slide 12 --- Key Technical Challenges \& How We Solved Them}
\begin{itemize}
\item \textbf{Temporal leakage} \emph{(most critical):} Features computed without a prediction date cutoff silently include future information --- metrics looked excellent, then collapsed on temporal validation; solution: strict \texttt{cutoff\_time} enforcement at every feature computation
\item \textbf{Patient-level data splitting:} Row-level random splits inflated metrics by allowing the model to "memorise" patients; solution: \texttt{GroupShuffleSplit} with patient ID as group key
\item \textbf{Missing data at scale:} \textasciitilde{}[X]\% of feature values are NaN; imputation introduced systematic bias; solution: XGBoost's native NaN handling --- no imputation needed
\item \textbf{OCR overcorrection:} The OCR fixer corrupted valid codes like "HbA1c" before a safelist was added; solution: \texttt{valid\_mixed\_terms} configuration parameter
\item \textbf{[Any challenges from Timarian's error analysis to add here]}
\end{itemize}

\noindent\rule{\linewidth}{0.4pt}

\subsection{Section 4 --- Results \& Evaluation (\textasciitilde{}11 min)}
\subsubsection{\textbf{\emph{ TRANSITION TO LIVE DASHBOARD HERE }}}
\emph{At this point, switch from the slide deck to the Streamlit monitoring dashboard. The following slide notes describe what to show on each dashboard tab --- they are speaker guide notes, not slides.}

\noindent\rule{\linewidth}{0.4pt}

\subsubsection{Dashboard Tab 1 --- Pipeline Overview / Model Performance}
\emph{Show the pipeline overview tab: patient count, observation count, extraction stats}
\begin{itemize}
\item Walk through the pipeline stage metrics: [X] patients loaded, [Y] clinical notes processed, [Z] vitals/labs extracted from notes
\item Transition to model performance tab
\item \textbf{Performance metrics framing:}
\begin{itemize}
\item Lead with the business metric: "For every 100 alerts at our operating threshold, [X] are genuine high-risk patients --- compared to [Y] for the HbA1c-only rule"
\item Precision at 80\% Recall: [result] --- the clinical threshold that balances catch rate with nursing workload
\item PR-AUC [result] vs. HbA1c baseline \textasciitilde{}0.155 --- [X]x improvement
\item ROC-AUC [result] --- reported for benchmark comparability
\end{itemize}
\item Show the baseline comparison table / chart
\end{itemize}

\noindent\rule{\linewidth}{0.4pt}

\subsubsection{Dashboard Tab 2 --- Patient Explorer / SHAP Demo \emph{(Alice \& Bob)}}
\emph{Show the patient explorer tab. Pull up a patient with a profile similar to Alice.}
\begin{itemize}
\item Select a HIGH or CRITICAL risk patient from the explorer
\item Show their timeline: HbA1c trend, encounter gaps, ER history
\item Show the SHAP force plot: "These are the three factors driving this alert --- HbA1c worsening over 180 days, 94 days since last encounter, eGFR declining"
\item Show the protective factors: "This is what is holding the risk below CRITICAL"
\item \textbf{Key message:} "This is what the clinician sees before deciding to act. Not a score. A reason."
\item \emph{"This is Alice. The system flagged her 9 days before her HbA1c crossed 10\%. Her care team was notified. She did not end up in hospital."}
\end{itemize}

\noindent\rule{\linewidth}{0.4pt}

\subsubsection{Dashboard Tab 3 --- Fairness, Robustness \& Compliance (if tab exists)}
\emph{Show subgroup performance or calibration plots if available on the dashboard}
\begin{itemize}
\item Demographic fairness: model performance across age groups, gender, ethnicity --- [insert finding]
\item Calibration: does the predicted 80\% risk correspond to an 80\% observed event rate?
\item Robustness: performance under degraded missing data rates --- which features are most critical?
\item \textbf{Compliance summary:} GDPR Art. 22 human review process; EU AI Act high-risk documentation; FHIR/SNOMED/ICD-10 compliance --- all evidenced in the report
\end{itemize}

\noindent\rule{\linewidth}{0.4pt}

\subsubsection{Back to Slides --- Error Analysis \& Limitations \emph{(brief, after dashboard)}}

\subsubsection{Slide 13 --- What the Model Gets Wrong}
\begin{itemize}
\item \textbf{False negatives (missed patients):} Sparse EHR data (low feature completeness), unusual comorbidity profiles, very recent first encounters --- patients the system has not had time to learn
\item \textbf{False positives (spurious alerts):} High medication count + stable HbA1c, elderly patients with high encounter frequency but no acute deterioration --- complex but stable cases
\item \textbf{Known limitations:}
\begin{itemize}
\item Trained on synthetic data --- real-world distributions will differ; positive rate in deployment may shift, requiring recalibration
\item Not validated for Type 1 diabetes, gestational diabetes, or paediatric patients
\item ZIP-code census features added no signal at 30-day horizon --- environmental factors are a long-horizon story, not a 30-day one
\end{itemize}
\item \textbf{[Insert Timarian's error analysis key findings]}
\end{itemize}

\noindent\rule{\linewidth}{0.4pt}

\subsection{Section 5 --- Key Takeaways (\textasciitilde{}6 min)}

\subsubsection{Slide 14 --- Our Recommendation, Restated \emph{(pyramid close)}}
\emph{Mirror Slide 2 --- close the loop for the client.}

\begin{mdframed}[style=recbox]
\textbf{"Deploy the 22-feature XGBoost pipeline. Pilot alert thresholds in one ward before hospital-wide rollout. Recalibrate if the real-world positive rate differs materially from the training distribution. The compliance documentation is ready for an EU AI Act audit."}
\end{mdframed}

\begin{itemize}
\item \textbf{ROI case:} [X] prevented readmissions per year $\times$ [CHF/USD Y average cost avoided] = [Z] net annual value; pipeline runs on existing hardware; no additional infrastructure required
\item \textbf{Risk:} Low --- system is a decision-support tool, not an autonomous decision-maker; all alerts require clinician review; no patient data leaves the hospital perimeter
\item \textbf{Next steps:} Ward pilot $\rightarrow$ threshold calibration $\rightarrow$ EHR integration via FHIR API $\rightarrow$ post-market monitoring setup
\end{itemize}

\noindent\rule{\linewidth}{0.4pt}

\subsubsection{Slide 15 --- Lessons for SwissCare \& Similar Engagements}
\begin{itemize}
\item \textbf{True positive rate over accuracy:} In clinical risk monitoring, missing a high-risk patient is categorically worse than a false alarm --- design metrics before models
\item \textbf{Explainability is a clinical requirement:} SHAP force plots are what make an alert actionable; a probability score alone is noise to a clinician under time pressure
\item \textbf{Temporal safety is not optional:} Features without a prediction-date cutoff produce circular reasoning that only fails when tested on the future; build it in from day one
\item \textbf{Synthetic data has limits:} Our census feature experiment proved that Synthea does not encode socioeconomic correlates of adverse events --- validate any external data source for signal before building on it
\item \textbf{Regulatory documentation is a first-class deliverable:} EU AI Act and GDPR requirements shape architecture (human oversight, logging, explainability) --- not just paperwork written after the fact
\end{itemize}

\noindent\rule{\linewidth}{0.4pt}

\subsection{Appendix / Backup Slides (for Q\&A)}

\begin{itemize}
\item Full 22-feature list with clinical meanings and SHAP global importance ranking
\item Detailed calibration curve
\item GDPR Article 22 human review process flowchart
\item Census feature correlation table (density\_per\_sq\_mile r=$-$0.0058, overall\_healthcare\_coverage r=+0.0050, coverage\_score r=+0.0034)
\item Failed experiment detail: SMOTE, row-level split, temporal cutoff omission, census features
\item Full regulatory compliance mapping (EU AI Act Annex III, GDPR Art. 22, FHIR schema, SNOMED event codes)
\item LSTM vs. XGBoost comparison
\item Stakeholder success metric mapping
\end{itemize}
\end{document}